\chapter{Testing and Verification}
\label{chap:testing}

This chapter documents how each component of this project was actually verified to behave correctly. All results shown reflect real test runs performed against the compiled binaries, not hypothetical output.

\section{Gopher Server: Manual Verification}

With the server running (\texttt{./gopher\_server gopher\_root} from within \texttt{gopher/}), the root menu was requested by sending an empty selector:

\begin{lstlisting}[style=bashstyle, numbers=none]
$ python3 -c "
import socket
s = socket.create_connection(('localhost', 7070))
s.sendall(b'\r\n')
print(s.recv(4096).decode())
"
1Welcome Menu	README	localhost	7070
0About this server	about.txt	localhost	7070
.
\end{lstlisting}

This confirms the menu format matches the specification described in Chapter~\ref{chap:gopher-protocol}: tab-separated fields, correct type characters, and the terminating \texttt{.} line.

Requesting an existing file and a nonexistent file were both verified:

\begin{lstlisting}[style=bashstyle, numbers=none]
$ ./gopher_client localhost 7070 about.txt
This is a sample Gopher file served over TCP.
.

$ ./gopher_client localhost 7070 nope.txt
3File not found
.
\end{lstlisting}

The server's own log confirmed each request was received and dispatched correctly:
\begin{lstlisting}[style=bashstyle, numbers=none]
[gopher] New connection from 127.0.0.1
[gopher] selector requested: "about.txt"
[gopher] New connection from 127.0.0.1
[gopher] selector requested: "nope.txt"
\end{lstlisting}

\section{Gopher Client Against the Gopher Server: End-to-End}

With the compiled \texttt{gopher\_client} (rather than an ad-hoc Python script) run against the compiled \texttt{gopher\_server}, all cases behaved identically to the manual tests above, and additionally the client's own argument-validation was exercised:

\begin{lstlisting}[style=bashstyle, numbers=none]
$ ./gopher_client
Usage: ./gopher_client <host> <port> [selector]
Example: ./gopher_client localhost 7070 about.txt
$ echo $?
1
\end{lstlisting}

This confirms the client correctly detects missing required arguments, prints a usage message to \texttt{stderr}, and exits with a nonzero status code --- the standard convention for signaling failure to calling shell scripts or automated tooling.

\section{HTTP Server: Status Codes and Headers}

With the server running (\texttt{./http\_server www} from within \texttt{http/}), each supported response path was exercised with \texttt{curl -i} (which prints response headers):

\begin{lstlisting}[style=bashstyle, numbers=none]
$ curl -s -i http://localhost:8080/
HTTP/1.1 200 OK
Content-Type: text/html
Content-Length: 175
Connection: close

<!DOCTYPE html>
...

$ curl -s -i http://localhost:8080/notfound.html
HTTP/1.1 404 Not Found
Content-Type: text/html
Content-Length: 48
Connection: close

<html><body><h1>404 Not Found</h1></body></html>

$ curl -s -i -X POST http://localhost:8080/
HTTP/1.1 405 Method Not Allowed
...
\end{lstlisting}

Each response's \inlinecode{Content-Length} header was confirmed to match the actual byte length of the body that followed, verifying the header-writing logic in \inlinecode{send\_response} (Chapter~\ref{chap:http-server}).

\section{The Path-Traversal Test: A Worked Example of a Misleading Result}

This test is documented in detail because the first attempt produced a result that looked like a failure but was not.

\textbf{First attempt}, using \texttt{curl}:
\begin{lstlisting}[style=bashstyle, numbers=none]
$ curl -s -i "http://localhost:8080/../etc/passwd"
HTTP/1.1 404 Not Found
...
\end{lstlisting}

At first glance, this looks wrong: the expectation was a \texttt{400 Bad Request} from the path-traversal check in \inlinecode{serve\_file} (Chapter~\ref{chap:http-server}). As explained in Issue 5 of Chapter~\ref{chap:running-on-linux}, the actual cause is that \texttt{curl} normalizes \texttt{..} segments out of the URL \emph{before} sending the request, so the server never actually receives a path containing \texttt{".."}.

\textbf{Second attempt}, bypassing \texttt{curl}'s normalization by sending a raw request directly over a socket:
\begin{lstlisting}[style=bashstyle, numbers=none]
$ python3 -c "
import socket
s = socket.create_connection(('localhost', 8080), timeout=3)
s.sendall(b'GET /../../etc/passwd HTTP/1.1\r\nHost: localhost\r\n\r\n')
print(s.recv(4096).decode())
"
HTTP/1.1 400 Bad Request
Content-Type: text/html
Content-Length: 50
Connection: close

<html><body><h1>400 Bad Request</h1></body></html>
\end{lstlisting}

This confirms the defense works correctly against a request that actually contains the literal \texttt{".."} substring the check is designed to catch. The server's log for this session shows both requests clearly, including the un-normalized path exactly as sent:

\begin{lstlisting}[style=bashstyle, numbers=none]
[http] GET /etc/passwd
[http] GET /../../etc/passwd
\end{lstlisting}

\begin{note}
This worked example is preserved here deliberately, rather than only showing the final passing test, because the debugging process --- recognizing that an apparently-failing test was actually a client-side artifact, and constructing a more direct test to isolate the real behavior --- is itself a core skill this project is meant to build.
\end{note}

\section{HTTP Client: Framing Verification}

The compiled \texttt{http\_client} was run against the compiled \texttt{http\_server} for each response type, with particular attention to the declared-vs-received length cross-check described in Chapter~\ref{chap:http-client}:

\begin{lstlisting}[style=bashstyle, numbers=none]
$ ./http_client localhost 8080 /
--- Headers ---
HTTP/1.1 200 OK
Content-Type: text/html
Content-Length: 175
Connection: close

--- Declared Content-Length: 175, actually received: 175 ---
--- Body ---
<!DOCTYPE html>
...

$ ./http_client localhost 8080 /about.txt
...
--- Declared Content-Length: 24, actually received: 24 ---
...

$ ./http_client localhost 8080 /notfound.html
...
--- Declared Content-Length: 48, actually received: 48 ---
...
\end{lstlisting}

In every case, the \texttt{Content-Length} value parsed from the response headers exactly matched the number of body bytes independently computed from the raw socket data (\inlinecode{body\_len\_from\_socket} in \texttt{http\_client.c}). This is the concrete, observable confirmation that the framing logic described conceptually in Chapter~\ref{chap:http-protocol} (Section~\ref{sec:http-framing}) and implemented in Chapters~\ref{chap:http-server} and~\ref{chap:http-client} is correct on both ends of the connection.

\section{Summary of Verified Behavior}

\begin{center}
\begin{tabularx}{\textwidth}{X l}
\toprule
\textbf{Behavior} & \textbf{Result} \\
\midrule
Gopher root menu format & Correct (tab-separated, correct types) \\
Gopher file retrieval (existing file) & Correct \\
Gopher file retrieval (missing file) & Correct (type-3 error) \\
Gopher client argument validation & Correct (usage message, exit code 1) \\
HTTP 200 response with correct headers & Correct \\
HTTP 404 for missing file & Correct \\
HTTP 405 for non-GET method & Correct \\
HTTP path-traversal defense (raw request) & Correct (400 Bad Request) \\
HTTP client Content-Length cross-check & Correct in all tested cases \\
\bottomrule
\end{tabularx}
\end{center}
